\chapter{Πρωτόκολλο μελέτης χρηστών}
\label{app:userstudy}

\section{Έντυπο συναίνεσης}
\todo{Από την επιτροπή ηθικής του ιδρύματος, αν απαιτείται.}

\section{Ερωτηματολόγιο SUS}

Κλίμακα 1 (διαφωνώ απόλυτα) έως 5 (συμφωνώ απόλυτα) \cite{brooke1996sus}.

\guide{Χρησιμοποίησε δημοσιευμένη, σταθμισμένη ελληνική μετάφραση του SUS αν υπάρχει, και αναφέρ' την. Αλλιώς η παρακάτω είναι δική σου μετάφραση και πρέπει να δηλωθεί ως τέτοια.}

\begin{enumerate}
  \item Νομίζω ότι θα ήθελα να χρησιμοποιώ αυτό το σύστημα συχνά.
  \item Βρήκα το σύστημα άσκοπα πολύπλοκο.
  \item Νομίζω ότι το σύστημα ήταν εύκολο στη χρήση.
  \item Νομίζω ότι θα χρειαζόμουν τη βοήθεια τεχνικού για να μπορώ να χρησιμοποιήσω αυτό το σύστημα.
  \item Βρήκα ότι οι διάφορες λειτουργίες του συστήματος ήταν καλά ενσωματωμένες.
  \item Νομίζω ότι υπήρχε υπερβολική ασυνέπεια στο σύστημα.
  \item Φαντάζομαι ότι οι περισσότεροι άνθρωποι θα μάθαιναν να χρησιμοποιούν αυτό το σύστημα πολύ γρήγορα.
  \item Βρήκα το σύστημα πολύ δύσχρηστο.
  \item Ένιωσα πολύ σίγουρος/η χρησιμοποιώντας το σύστημα.
  \item Χρειάστηκε να μάθω πολλά πράγματα πριν μπορέσω να ξεκινήσω με αυτό το σύστημα.
\end{enumerate}

Βαθμολογία: για τα περιττά θέματα $(x-1)$, για τα άρτια $(5-x)$· το άθροισμα επί $2{,}5$ δίνει τιμή $0$--$100$.

\section{NASA-TLX}
\guide{Έξι κλίμακες 0--100: νοητικός φόρτος, σωματικός φόρτος, χρονική πίεση, απόδοση, προσπάθεια, απογοήτευση \cite{hart1988nasatlx}. Χρησιμοποίησε την απλή μορφή (raw TLX, μέσος όρος) εκτός αν ο επιβλέπων ζητήσει τις σταθμίσεις.}

\section{Εργασίες}
\begin{enumerate}
  \item Βρες τους αισθητήρες \fillin{Κ3}, \fillin{Κ7}, \fillin{Κ12} και ακούμπησέ τους.
  \item Ποιοι αισθητήρες είναι σε κατάσταση συναγερμού και πού βρίσκονται;
  \item Ποια περιοχή είχε τη μεγαλύτερη μετακίνηση τον τελευταίο μήνα;
\end{enumerate}

\section{Ημιδομημένη συνέντευξη}
\todo{5--6 ανοιχτές ερωτήσεις: τι βοήθησε, τι δυσκόλεψε, θα το χρησιμοποιούσες στη δουλειά, τι λείπει.}
